\section{Conclusion}
\label{sec:conclusion}

This review connected general generation approaches to automotive LiDAR
and its evaluation. Simulation, learned generators and hybrids offer
different balances of scene control, coverage and dependence on real data.
Useful training datasets require sensor sampling and annotation consistency
in addition to plausible geometry.

The selected studies support synthetic pretraining and supplementation in
particular perception settings, but do not establish a universal ranking
or replacement for real data. Similarity and completion evaluations provide
complementary evidence, while controlled assessment on independent real
observations most directly establishes training utility. Connecting sensor
diagnostics and scenario coverage to consistent downstream gains could
help identify which aspects of synthesis merit improvement and where real
measurements remain necessary.
